%% ma: L11-B04-0009
%% lop: 11
%% chuong: 1
%% bai: 4
%% dang: 11.4.9
%% loai: DS
%% muc: [TH, TH, VD, TH]
%% dap_an: "SSĐĐ"
%% nguon: "de-cuong-giua-ky-1-toan-11-nam-2026-2027-truong-thpt-hoang-van-thu-ha-noi-2-8.pdf (M08, Đề cương giữa kì 1 Toán 11 THPT Hoàng Văn Thụ, Hà Nội), Phần 2 Câu 4"
%% kien_thuc: "Bài 4 (phương trình lượng giác cơ bản)"
%% trang_thai: chinh-thuc
%% ghi_chu: "Đã duyệt 29/10/2026."
\begin{ex}
Cho phương trình lượng giác $2\cos x=\sqrt3$, khi đó:
\choiceTF{Phương trình có nghiệm $x=\pm\dfrac{\pi}{3}+k2\pi$ $(k\in\mathbb Z)$}{Trong đoạn $\left[0;\dfrac{5\pi}{2}\right]$ phương trình có $4$ nghiệm}{\True Tổng các nghiệm của phương trình trong đoạn $\left[0;\dfrac{5\pi}{2}\right]$ bằng $\dfrac{25\pi}{6}$}{\True Trong đoạn $\left[0;\dfrac{5\pi}{2}\right]$ phương trình có nghiệm lớn nhất bằng $\dfrac{13\pi}{6}$}
\loigiai{a) Sai: $\cos x=\dfrac{\sqrt3}2\Leftrightarrow x=\pm\dfrac\pi6+k2\pi$ $(k\in\mathbb Z)$.

b) Sai: trên $\left[0;\dfrac{5\pi}2\right]$ các nghiệm là $\dfrac\pi6$, $\dfrac{11\pi}6$, $\dfrac{13\pi}6$; có $3$ nghiệm.

c) Đúng: $\dfrac\pi6+\dfrac{11\pi}6+\dfrac{13\pi}6=\dfrac{25\pi}6$.

d) Đúng: nghiệm lớn nhất là $\dfrac{13\pi}{6}$.}
\end{ex}
